\documentclass[11pt]{article}
\usepackage[T1]{fontenc}
\usepackage{textcomp}
\usepackage[margin=0.75in]{geometry}
\usepackage{amsmath,amssymb,amsthm}
\usepackage{array,booktabs}
\usepackage{hyperref}
\setlength{\parindent}{0pt}
\newtheorem{theorem}{Theorem}
\theoremstyle{definition}
\newtheorem{definition}{Definition}
\title{Flow Proof Artifact: Eq-reflexive-one}
\author{Generated by \texttt{flow doc proof}}
\date{Flow Proof Book}
\begin{document}
\maketitle
One equals itself.\\[0.5em]
\textbf{Source.} leibniz — https://en.wikipedia.org/wiki/Law\_of\_identity\\[0.5em]
\bigskip
\noindent{\large \textbf{Derived fact 1.} \textit{1 = 1} \hfill equality \cdot equality \cdot one equals itself}\par
\smallskip\noindent\textit{Coordinate: equality \cdot equality \cdot one equals itself}\par
\smallskip\noindent\textit{Source: leibniz}\par
\smallskip\noindent\textit{Built on: anything is always equal to itself}\par
\medskip
\noindent\textbf{Goal.} 1 = 1\par
\noindent$\displaystyle \mathrm{succ}(0) = \mathrm{succ}(0)$\par
\medskip
\noindent
\begin{tabular}{@{} >{\raggedright\arraybackslash}p{0.06\textwidth} >{\raggedright\arraybackslash}p{0.47\textwidth} >{\raggedleft\arraybackslash}p{0.41\textwidth} @{}}
\textbf{\#} & \textbf{Proof} & \textbf{Mathematics} \\
\midrule
\textbf{1.} & We prove that one equals itself for equality on equality. & \\[0.45em]
\textbf{2.} & We invoke the axiom governing equality on equality: anything is always equal to itself (instantiated for succ(0)). & $\displaystyle \mathrm{succ}(0) = \mathrm{succ}(0)$ \\[0.45em]
\textbf{3.} & From step 2, this implies the successor of 0 equals the successor of 0. Hence proven. & $\displaystyle \mathrm{succ}(0) = \mathrm{succ}(0)$ \\[0.45em]
\bottomrule
\end{tabular}
\label{thm:Eq-equality-one_equals_itself}
\par\medskip\hrule
\end{document}
